\documentclass[aspectratio=169, xcolor=dvipsnames]{beamer}

% --- Theme Setup ---
\usetheme{Madrid}
\usecolortheme{default}

% --- Packages ---
\usepackage{graphicx}
\usepackage{xcolor}
\usepackage{booktabs}
\usepackage{hyperref}
\usepackage{adjustbox}
\usepackage{amssymb}
\usepackage{amsmath}
\usepackage{listings}
\usepackage{caption}
\usepackage{tikz}

\lstset{
  basicstyle=\ttfamily\small,
  keywordstyle=\color{blue}\bfseries,
  stringstyle=\color{red},
  showstringspaces=false,
  breaklines=true
}

% --- Title Information ---
\title{Module 48}
\subtitle{Transactions/3: Recoverability}
\author{Partha Pratim Das}
\institute{Department of Computer Science and Engineering\\Indian Institute of Technology, Kharagpur\\ppd@cse.iitkgp.ac.in}
\date{\today}

\begin{document}

% Title Slide
\begin{frame}
    \titlepage
\end{frame}

\begin{frame}{Module Recap}
    \begin{itemize}[<+->]
        \item Understood the issues that arise when two or more transactions work concurrently
        \item Learnt the forms of serializability in terms of conflict and view serializability
        \item Acyclic precedence graph can ensure conflict serializability
    \end{itemize}
\end{frame}

\begin{frame}{Module Objectives}
    \begin{itemize}[<+->]
        \item What happens if system fails while a transaction is in execution? Can a consistent state be reached for the database? Recoverability attempts to answer issues in state and transaction recovery in the face of system failures
        \item Conflict serializability is a crisp concept for concurrent execution that guarantees ACID properties and has a simple detection algorithm. Yet only few schedules are Conflict serializable in practice. There is a need to explore - View Serializability - a weaker system for better concurrency
    \end{itemize}
\end{frame}

\begin{frame}{Module Outline}
    \begin{itemize}[<+->]
        \item Recoverability
        \item Transaction Definition in SQL
        \item View Serializability
        \item Complex Notions of Serializability
    \end{itemize}
\end{frame}

\section{Recovery}

\begin{frame}
  \vfill
  \centering
  \begin{beamercolorbox}[sep=8pt,center,shadow=true,rounded=true]{title}
    \usebeamerfont{title}Recovery\par%
  \end{beamercolorbox}
  \vfill
\end{frame}

\begin{frame}{What is Recovery?}
    \begin{itemize}[<+->]
        \item Serializability helps to ensure Isolation and Consistency of a schedule
        \item Yet, the Atomicity and Consistency may be compromised in the face of system failures
        \item Consider a schedule comprising a single transaction (obviously serial):
        \begin{enumerate}
            \item \texttt{read(A)}
            \item \texttt{A := A - 50}
            \item \texttt{write(A)}
            \item \texttt{read(B)}
            \item \texttt{B := B + 50}
            \item \texttt{write(B)}
            \item \texttt{commit} // Make the changes permanent; show the results to the user
        \end{enumerate}
        \item What if system fails after Step 3 and before Step 6?
        \item Leads to inconsistent state
        \item Need to rollback update of A
        \item This is known as Recovery
    \end{itemize}
\end{frame}

\begin{frame}{Recoverable Schedules}
    \begin{itemize}[<+->]
        \item If a transaction $T_j$ reads a data item previously written by a transaction $T_i$, then the commit operation of $T_i$ must appear before the commit operation of $T_j$.
        \item The following schedule is not recoverable if $T_9$ commits immediately after the \texttt{read(A)} operation:
    \end{itemize}
    \vspace{0.5em}
    \uncover<3->{
        \begin{center}
            \begin{tabular}{|l|l|}
            \hline
            \textbf{$T_8$} & \textbf{$T_9$} \\ \hline
            \texttt{read(A)} & \\ 
            \texttt{write(A)} & \\ 
            & \texttt{read(A)} \\ 
            & \texttt{commit} \\ 
            \texttt{read(B)} & \\ \hline
            \end{tabular}
        \end{center}
    }
    \begin{itemize}[<+->]
        \item If $T_8$ should abort, $T_9$ would have read (and possibly shown to the user) an inconsistent database state. Hence, database must ensure that schedules are recoverable.
    \end{itemize}
\end{frame}

\begin{frame}{Cascading Rollbacks}
    \begin{block}{Cascading Rollback}
        A single transaction failure leads to a series of transaction rollbacks.
    \end{block}
    \begin{itemize}[<+->]
        \item Consider the following schedule where none of the transactions has yet committed (so the schedule is recoverable):
    \end{itemize}
    \uncover<3->{
        \begin{center}
            \begin{tabular}{|l|l|l|}
            \hline
            \textbf{$T_{10}$} & \textbf{$T_{11}$} & \textbf{$T_{12}$} \\ \hline
            \texttt{read(A)} & & \\ 
            \texttt{read(B)} & & \\ 
            \texttt{write(A)} & & \\ 
            & \texttt{read(A)} & \\ 
            & \texttt{write(A)} & \\ 
            & & \texttt{read(A)} \\ 
            \texttt{abort} & & \\ \hline
            \end{tabular}
        \end{center}
    }
    \begin{itemize}[<+->]
        \item If $T_{10}$ fails, $T_{11}$ and $T_{12}$ must also be rolled back
        \item Can lead to the undoing of a significant amount of work
    \end{itemize}
\end{frame}

\begin{frame}{Cascadeless Schedules}
    \begin{block}{Cascadeless Schedules}
        For each pair of transactions $T_i$ and $T_j$ such that $T_j$ reads a data item previously written by $T_i$, the commit operation of $T_i$ appears before the read operation of $T_j$.
    \end{block}
    \begin{itemize}[<+->]
        \item Every cascadeless schedule is also recoverable
        \item It is desirable to restrict the schedules to those that are cascadeless
        \item Example of a schedule that is NOT cascadeless (same as previous slide):
    \end{itemize}
    \uncover<4->{
        \begin{center}
            \begin{tabular}{|l|l|l|}
            \hline
            \textbf{$T_{10}$} & \textbf{$T_{11}$} & \textbf{$T_{12}$} \\ \hline
            \texttt{read(A)} &  & \\ 
            \texttt{read(B)} &  & \\ 
            \texttt{write(A)} &  & \\ 
            \texttt{abort} & \texttt{read(A)} & \\ 
            & \texttt{write(A)} & \texttt{read(A)} \\ \hline
            \end{tabular}
        \end{center}
    }
\end{frame}

\begin{frame}{Example: Irrecoverable Schedule}
    \begin{center}
        \begin{adjustbox}{max width=0.9\textwidth}
        \begin{tabular}{|l|l|l|l|}
        \hline
        \textbf{T1} & \textbf{T2} & \textbf{T2's Buffer} & \textbf{Database} \\ \hline
        & & & A = 5000 \\ \hline
        \texttt{R(A);} & A = 5000 & & A = 5000 \\ \hline
        \texttt{A = A - 1000;} & A = 4000 & & A = 5000 \\ \hline
        \texttt{W(A);} & A = 4000 & & A = 4000 \\ \hline
        & \texttt{R(A);} & A = 4000 & A = 4000 \\ \hline
        & & A = 4500 & A = 4000 \\ \hline
        & \texttt{W(A);} & A = 4500 & A = 4500 \\ \hline
        & \texttt{Commit;} & & \\ \hline
        Failure Point & & & \\ \hline
        \texttt{Commit;} & & & \\ \hline
        \end{tabular}
        \end{adjustbox}
    \end{center}
    \vspace{1em}
    \uncover<2->{
        Rollback is possible only till the end (commit) of T2. So the computation of A (4000) and write in T1 is lost.
    }
\end{frame}

\begin{frame}{Example: Recoverable Schedule with Cascading Rollback}
    \begin{center}
        \begin{adjustbox}{max width=0.9\textwidth}
        \begin{tabular}{|l|l|l|l|}
        \hline
        \textbf{T1} & \textbf{T2} & \textbf{T2's Buffer} & \textbf{Database} \\ \hline
        & & & A = 5000 \\ \hline
        \texttt{R(A);} & A = 5000 & & A = 5000 \\ \hline
        \texttt{A = A - 1000;} & A = 4000 & & A = 5000 \\ \hline
        \texttt{W(A);} & A = 4000 & & A = 4000 \\ \hline
        & \texttt{R(A);} & A = 4000 & A = 4000 \\ \hline
        & & A = 4500 & A = 4000 \\ \hline
        & \texttt{W(A);} & A = 4500 & A = 4500 \\ \hline
        Failure Point & & & \\ \hline
        \texttt{Commit;} & & & \\ \hline
        & \texttt{Commit;} & & \\ \hline
        \end{tabular}
        \end{adjustbox}
    \end{center}
    \vspace{1em}
    \uncover<2->{
        Rollback is possible as T2 has not committed yet. But T2 also needs to be rolled back for rolling back T1.
    }
\end{frame}

\begin{frame}{Example: Recoverable Schedule without Cascading Rollback}
    \begin{center}
        \begin{adjustbox}{max width=0.9\textwidth}
        \begin{tabular}{|l|l|l|l|}
        \hline
        \textbf{T1} & \textbf{T2} & \textbf{T2's Buffer} & \textbf{Database} \\ \hline
        & & & A = 5000 \\ \hline
        \texttt{R(A);} & & & A = 5000 \\ \hline
        \texttt{A = 4000;} & & & A = 5000 \\ \hline
        \texttt{W(A);} & & & A = 4000 \\ \hline
        \texttt{Commit;} & & & \\ \hline
        & \texttt{R(A);} & A = 4000 & A = 4000 \\ \hline
        & \texttt{A = A + 500;} & A = 4500 & A = 4000 \\ \hline
        & \texttt{W(A);} & A = 4500 & A = 4500 \\ \hline
        & \texttt{Commit;} & & \\ \hline
        \end{tabular}
        \end{adjustbox}
    \end{center}
    \vspace{1em}
    \uncover<2->{
        Rollback is possible without cascading - wherever failure occurs.
    }
\end{frame}

\section{Transaction Definition in SQL}

\begin{frame}
  \vfill
  \centering
  \begin{beamercolorbox}[sep=8pt,center,shadow=true,rounded=true]{title}
    \usebeamerfont{title}Transaction Definition in SQL\par%
  \end{beamercolorbox}
  \vfill
\end{frame}

\begin{frame}{Transaction Definition in SQL}
    \begin{itemize}[<+->]
        \item Data manipulation language must include a construct for specifying the set of actions that comprise a transaction
        \item In SQL, a transaction begins implicitly
        \item A transaction in SQL ends by:
        \begin{itemize}
            \item \textbf{Commit work}
            \begin{itemize}
                \item[$\triangleright$] Commits the current transaction and begins a new one
            \end{itemize}
            \item \textbf{Rollback work}
            \begin{itemize}
                \item[$\triangleright$] Causes current transaction to abort
            \end{itemize}
        \end{itemize}
        \item In almost all database systems, by default, every SQL statement also commits implicitly if it executes successfully
        \begin{itemize}
            \item[$\triangleright$] Implicit commit can be turned off by a database directive
            \item[$\triangleright$] For example in JDBC, \texttt{connection.setAutoCommit(false);}
        \end{itemize}
    \end{itemize}
\end{frame}

\begin{frame}{Transaction Control Language (TCL)}
    \begin{itemize}[<+->]
        \item The following commands are used to control transactions
        \begin{itemize}
            \item \textbf{COMMIT}
            \begin{itemize}
                \item[$\triangleright$] To save the changes
            \end{itemize}
            \item \textbf{ROLLBACK}
            \begin{itemize}
                \item[$\triangleright$] To roll back the changes
            \end{itemize}
            \item \textbf{SAVEPOINT}
            \begin{itemize}
                \item[$\triangleright$] Creates points within the groups of transactions in which to ROLLBACK
            \end{itemize}
            \item \textbf{SET TRANSACTION}
            \begin{itemize}
                \item[$\triangleright$] Places a name on a transaction
            \end{itemize}
        \end{itemize}
        \item Transactional control commands are only used with the DML Commands such as \texttt{INSERT}, \texttt{UPDATE} and \texttt{DELETE} only
        \item They cannot be used while creating tables or dropping them because these operations are automatically committed in the database
    \end{itemize}
\end{frame}

\begin{frame}[fragile]{TCL: COMMIT Command}
    \begin{itemize}[<+->]
        \item \texttt{COMMIT} is the transactional command used to save changes invoked by a transaction to the database
        \item \texttt{COMMIT} saves all the transactions to the database since the last \texttt{COMMIT} or \texttt{ROLLBACK} command
        \item The syntax for the \texttt{COMMIT} command is as follows:
        \begin{itemize}
            \item \texttt{SQL> DELETE FROM Customers WHERE AGE = 25;}
            \item \texttt{SQL> COMMIT;}
        \end{itemize}
    \end{itemize}
    \vspace{0.5em}
    \uncover<5->{
        \textbf{\texttt{SQL> SELECT * FROM Customers;}}
        \begin{center}
            \begin{tabular}{|l|l|l|l|l|}
            \hline
            \textbf{ID} & \textbf{NAME} & \textbf{AGE} & \textbf{ADDRESS} & \textbf{SALARY} \\ \hline
            1 & Ramesh & 32 & Ahmedabad & 2000 \\ \hline
            3 & kaushik & 23 & Kota & 2000 \\ \hline
            5 & Hardik & 27 & Bhopal & 8500 \\ \hline
            6 & Komal & 22 & MP & 4500 \\ \hline
            7 & Muffy & 24 & Indore & 10000 \\ \hline
            \end{tabular}
        \end{center}
    }
\end{frame}

\begin{frame}[fragile]{TCL: ROLLBACK Command}
    \begin{itemize}[<+->]
        \item The \texttt{ROLLBACK} is the command used to undo transactions that have not already been saved to the database
        \item This can only be used to undo transactions since the last \texttt{COMMIT} or \texttt{ROLLBACK} command was issued
        \item The syntax for a \texttt{ROLLBACK} command is as follows:
        \begin{itemize}
            \item \texttt{SQL> DELETE FROM Customers WHERE AGE = 25;}
            \item \texttt{SQL> ROLLBACK;}
        \end{itemize}
    \end{itemize}
    \vspace{0.5em}
    \uncover<5->{
        \textbf{\texttt{SQL> SELECT * FROM Customers;}}
        \begin{center}
            \begin{tabular}{|l|l|l|l|l|}
            \hline
            \textbf{ID} & \textbf{NAME} & \textbf{AGE} & \textbf{ADDRESS} & \textbf{SALARY} \\ \hline
            1 & Ramesh & 32 & Ahmedabad & 2000 \\ \hline
            2 & Khilan & 25 & Delhi & 1500 \\ \hline
            3 & kaushik & 23 & Kota & 2000 \\ \hline
            4 & Chaitali & 25 & Mumbai & 6500 \\ \hline
            5 & Hardik & 27 & Bhopal & 8500 \\ \hline
            6 & Komal & 22 & MP & 4500 \\ \hline
            7 & Muffy & 24 & Indore & 10000 \\ \hline
            \end{tabular}
        \end{center}
    }
\end{frame}

\begin{frame}[fragile]{TCL: SAVEPOINT / ROLLBACK Command}
    \begin{itemize}[<+->]
        \item A \texttt{SAVEPOINT} is a point in a transaction when you can roll the transaction back to a certain point without rolling back the entire transaction
        \item The syntax for a \texttt{SAVEPOINT} command is:
        \begin{itemize}
            \item \texttt{SAVEPOINT SAVEPOINT\_NAME;}
        \end{itemize}
        \item This command serves only in the creation of a SAVEPOINT among all the transactional statements.
        \item The \texttt{ROLLBACK} command is used to undo a group of transactions
        \item The syntax for rolling back to a \texttt{SAVEPOINT} is:
        \begin{itemize}
            \item \texttt{ROLLBACK TO SAVEPOINT\_NAME;}
        \end{itemize}
    \end{itemize}
\end{frame}

\begin{frame}[fragile]{TCL: SAVEPOINT / ROLLBACK Command Example}
    \begin{exampleblock}{Example}
\begin{verbatim}
SQL> SAVEPOINT SP1;
Savepoint created.
SQL> DELETE FROM Customers WHERE ID=1;
1 row deleted.
SQL> SAVEPOINT SP2;
Savepoint created.
SQL> DELETE FROM Customers WHERE ID=2;
1 row deleted.
SQL> SAVEPOINT SP3;
Savepoint created.
SQL> DELETE FROM Customers WHERE ID=3;
1 row deleted.
\end{verbatim}
    \end{exampleblock}
\end{frame}

\begin{frame}[fragile]{TCL: SAVEPOINT / ROLLBACK Command Example (cont.)}
    \begin{itemize}[<+->]
        \item Three records deleted
        \item Undo the deletion of last two:
        \begin{itemize}
            \item \texttt{SQL> ROLLBACK TO SP2;}
            \item \texttt{Rollback complete}
        \end{itemize}
    \end{itemize}
    \vspace{0.5em}
    \uncover<4->{
        \textbf{\texttt{SQL> SELECT * FROM Customers;}}
        \begin{center}
            \begin{tabular}{|l|l|l|l|l|}
            \hline
            \textbf{ID} & \textbf{NAME} & \textbf{AGE} & \textbf{ADDRESS} & \textbf{SALARY} \\ \hline
            2 & Khilan & 25 & Delhi & 1500 \\ \hline
            3 & kaushik & 23 & Kota & 2000 \\ \hline
            4 & Chaitali & 25 & Mumbai & 6500 \\ \hline
            5 & Hardik & 27 & Bhopal & 8500 \\ \hline
            6 & Komal & 22 & MP & 4500 \\ \hline
            7 & Muffy & 24 & Indore & 10000 \\ \hline
            \end{tabular}
        \end{center}
    }
\end{frame}

\begin{frame}{TCL: RELEASE SAVEPOINT Command}
    \begin{itemize}[<+->]
        \item The \texttt{RELEASE SAVEPOINT} command is used to remove a \texttt{SAVEPOINT} that you have created
        \item The syntax for a \texttt{RELEASE SAVEPOINT} command is as follows:
        \begin{itemize}
            \item \texttt{RELEASE SAVEPOINT SAVEPOINT\_NAME;}
        \end{itemize}
        \item Once a \texttt{SAVEPOINT} has been released, you can no longer use the \texttt{ROLLBACK} command to undo transactions performed since the last \texttt{SAVEPOINT}
    \end{itemize}
\end{frame}

\begin{frame}{TCL: SET TRANSACTION Command}
    \begin{itemize}[<+->]
        \item The \texttt{SET TRANSACTION} command can be used to initiate a database transaction
        \item This command is used to specify characteristics for the transaction that follows
        \item For example, you can specify a transaction to be read only or read write
        \item The syntax for a \texttt{SET TRANSACTION} command is as follows:
        \begin{itemize}
            \item \texttt{SET TRANSACTION [ READ WRITE | READ ONLY ];}
        \end{itemize}
    \end{itemize}
\end{frame}

\section{View Serializability}

\begin{frame}
  \vfill
  \centering
  \begin{beamercolorbox}[sep=8pt,center,shadow=true,rounded=true]{title}
    \usebeamerfont{title}View Serializability\par%
  \end{beamercolorbox}
  \vfill
\end{frame}

\begin{frame}{View Serializability}
    \begin{itemize}[<+->]
        \item Let S and S$'$ be two schedules with the same set of transactions. S and S$'$ are \textbf{view equivalent} if the following three conditions are met, for each data item Q:
        \begin{itemize}
            \item[$\triangleright$] \textbf{Initial Read}: If in schedule S, transaction $T_i$ reads the initial value of Q, then in schedule S$'$ also transaction $T_i$ must read the initial value of Q
            \item[$\triangleright$] \textbf{Write-Read Pair}: If in schedule S transaction $T_i$ executes \texttt{read(Q)}, and that value was produced by transaction $T_j$ (if any), then in schedule S$'$ also transaction $T_i$ must read the value of Q that was produced by the same \texttt{write(Q)} operation of transaction $T_j$
            \item[$\triangleright$] \textbf{Final Write}: The transaction (if any) that performs the final \texttt{write(Q)} operation in schedule S must also perform the final \texttt{write(Q)} operation in schedule S$'$
        \end{itemize}
        \item As can be seen, view equivalence is also based purely on reads and writes alone
    \end{itemize}
\end{frame}

\begin{frame}{View Serializability (2)}
    \begin{itemize}[<+->]
        \item A schedule S is \textbf{view serializable} if it is view equivalent to a serial schedule
        \item Every conflict serializable schedule is also view serializable
        \item Below is a schedule which is view-serializable but not conflict serializable:
    \end{itemize}
    \uncover<4->{
        \begin{center}
            \begin{tabular}{|l|l|l|}
            \hline
            \textbf{$T_{27}$} & \textbf{$T_{28}$} & \textbf{$T_{29}$} \\ \hline
            \texttt{read(Q)} & & \\ \hline
            \texttt{write(Q)} & & \\ \hline
            & \texttt{write(Q)} & \\ \hline
            & & \texttt{write(Q)} \\ \hline
            \end{tabular}
        \end{center}
    }
    \begin{itemize}[<+->]
        \item What serial schedule is above equivalent to?
        \begin{itemize}
            \item[$\triangleright$] $T_{27} \rightarrow T_{28} \rightarrow T_{29}$
        \end{itemize}
        \item The one \texttt{read(Q)} instruction reads the initial value of Q in both schedules and $T_{29}$ performs the final write of Q in both schedules
        \item $T_{28}$ and $T_{29}$ perform \texttt{write(Q)} operations called \textbf{blind writes}, without having performed a \texttt{read(Q)} operation
        \item Every view serializable schedule that is not conflict serializable has blind writes
    \end{itemize}
\end{frame}

\begin{frame}{Test for View Serializability}
    \begin{itemize}[<+->]
        \item The precedence graph test for conflict serializability cannot be used directly to test for view serializability
        \item Extension to test for view serializability has cost exponential in the size of the precedence graph
        \item The problem of checking if a schedule is view serializable falls in the class of NP-complete problems
        \item Thus, existence of an efficient algorithm is extremely unlikely
        \item However, practical algorithms that just check some sufficient conditions for view serializability can still be used
    \end{itemize}
\end{frame}

\begin{frame}{View Serializability: Example 1}
    \begin{exampleblock}{Example}
        \begin{itemize}[<+->]
            \item Check whether the schedule is view serializable or not?
            \item S : \texttt{R2(B); R2(A); R1(A); R3(A); W1(B); W2(B); W3(B);}
        \end{itemize}
    \end{exampleblock}
    \uncover<3->{
        \textbf{Solution:}
        \begin{itemize}[<+->]
            \item With 3 transactions, total number of schedules possible = 3! = 6
            \begin{itemize}
                \item[$\triangleright$] $< T_1 T_2 T_3 >$
                \item[$\triangleright$] $< T_1 T_3 T_2 >$
                \item[$\triangleright$] $< T_2 T_3 T_1 >$
                \item[$\triangleright$] $< T_2 T_1 T_3 >$
                \item[$\triangleright$] $< T_3 T_1 T_2 >$
                \item[$\triangleright$] $< T_3 T_2 T_1 >$
            \end{itemize}
        \end{itemize}
    }
\end{frame}

\begin{frame}{View Serializability: Example 1 (2)}
    \textbf{Solution:}
    \begin{itemize}[<+->]
        \item Final update on data items:
        \begin{itemize}
            \item[$\triangleright$] A : (No write on A)
            \item[$\triangleright$] B : $T_1$, $T_2$, $T_3$ (All 3 transactions write B)
            \item[$\triangleright$] As the final update on B is made by $T_3$, ($T_1$, $T_2$) $\rightarrow$ $T_3$. Now, Removing those schedules in which $T_3$ is not executing at last:
            \begin{itemize}
                \item $< T_1 T_2 T_3 >$
                \item $< T_2 T_1 T_3 >$
            \end{itemize}
        \end{itemize}
    \end{itemize}
\end{frame}

\begin{frame}{View Serializability: Example 1 (3)}
    \textbf{Solution:}
    \begin{itemize}[<+->]
        \item Initial Read + Which transaction updates after read?
        \begin{itemize}
            \item[$\triangleright$] A : $T_2$, $T_1$, $T_3$ (initial read)
            \item[$\triangleright$] B : $T_2$ (initial read); $T_1$ (update after read)
            \item[$\triangleright$] The transaction $T_2$ reads B initially which is updated by $T_1$. So $T_2$ must execute before $T_1$. Hence, $T_2 \rightarrow T_1$. So only one schedule survives:
            \begin{itemize}
                \item $< T_2 T_1 T_3 >$
            \end{itemize}
        \end{itemize}
        \item Write Read Sequence (WR)
        \begin{itemize}
            \item[$\triangleright$] No need to check here
        \end{itemize}
        \item Hence, view equivalent serial schedule is:
        \begin{itemize}
            \item[$\triangleright$] $T_2 \rightarrow T_1 \rightarrow T_3$
        \end{itemize}
    \end{itemize}
\end{frame}

\begin{frame}{View Serializability: Example 2}
    \begin{exampleblock}{Example}
        \begin{itemize}[<+->]
            \item Check whether S is conflict serializable and/or view serializable or not?
            \item S : \texttt{R1(A); R2(A); R3(A); R4(A); W1(B); W2(B); W3(B); W4(B)}
            \item Solution is given in the next slide (hidden). First try to solve this and then check the solution.
        \end{itemize}
    \end{exampleblock}
\end{frame}

\section{Complex Notions of Serializability}

\begin{frame}
  \vfill
  \centering
  \begin{beamercolorbox}[sep=8pt,center,shadow=true,rounded=true]{title}
    \usebeamerfont{title}Complex Notions of Serializability\par%
  \end{beamercolorbox}
  \vfill
\end{frame}

\begin{frame}{More Complex Notions of Serializability}
    \begin{itemize}[<+->]
        \item The schedule below produces the same outcome as the serial schedule $< T_1, T_5 >$, yet is not conflict equivalent or view equivalent to it
        \item If we start with A = 1000 and B = 2000, the final result is 960 and 2040
        \item Determining such equivalence requires analysis of operations other than read and write
    \end{itemize}
    \uncover<4->{
        \begin{center}
            \begin{tabular}{|l|l|}
            \hline
            \textbf{$T_1$} & \textbf{$T_5$} \\ \hline
            \texttt{read(A)} & \\ 
            \texttt{write(A)} & \\ 
            & \texttt{read(B)} \\ 
            & \texttt{B := B - 10} \\ 
            & \texttt{write(B)} \\ 
            \texttt{read(B)} & \\ 
            \texttt{B := B + 50} & \\ 
            \texttt{write(B)} & \\ 
            & \texttt{read(A)} \\ 
            & \texttt{A := A + 10} \\ 
            & \texttt{write(A)} \\ \hline
            \end{tabular}
        \end{center}
    }
\end{frame}

\begin{frame}{Module Summary}
    \begin{itemize}[<+->]
        \item With proper planning, a database can be recovered back to a consistent state from inconsistent state in the face of system failures. Such a recovery is done via cascaded or cascadeless rollback
        \item View Serializability is a weaker serializability system for better concurrency. However, testing for view serializability is NP complete
    \end{itemize}
    \vspace{2em}
    \uncover<3->{
        \begin{center}
        \small \textit{Slides used in this presentation are borrowed from \url{http://db-book.com/} with kind permission of the authors.} \\
        \small \textit{Edited and new slides are marked with 'PPD'.}
        \end{center}
    }
\end{frame}

\end{document}
